\documentclass{article}
\usepackage[T1]{fontenc}
\usepackage{lmodern}
\usepackage{graphicx}
\usepackage{amsmath,amssymb}
\usepackage[a4paper,margin=2cm]{geometry}
\usepackage[absolute,overlay]{textpos}
\setlength{\TPHorizModule}{1mm}
\setlength{\TPVertModule}{1mm}
\usepackage[indonesian]{babel}
\usepackage{hyperref}
\setlength{\emergencystretch}{2em}
% unit: Halaman 36

\begin{document}

% === Halaman 36 (python/native) ===

36

Kunci-kunci lemah DES

• Kunci lemah DES adalah kunci K sedemikian sehingga EK(EK(x)) = x untuk semua 

plainteks x, yang artinya enkripsi dan dekripsi menjadi sama.

• Kunci lemah ini menghasilkan pembangkitan kunci putaran yang sama  pada 

semua putaran.

• DES memiliki 4 kunci lemah (masing-masing panjangnya 56 bit), dinyatakan dalam 

kode heksadesimal sbb:


00000000000000 


0000000FFFFFFF 


FFFFFFF0000000 


FFFFFFFFFFFFFF 

• Kunci lemah seharusnya dihindari.

\end{document}
